\chapter{Tinnitus (Ringing in the Ears)}

\section*{Definition}
Perception of sound in the absence of an external acoustic stimulus, described as ringing, buzzing, hissing, humming, or clicking, which may be unilateral or bilateral, continuous or intermittent.

\section*{Pathophysiology}
Abnormal neural activity in the auditory pathway, from cochlear hair cell damage to central auditory cortex reorganization. Subtypes include objective (arising from vascular or muscular sources near the ear) and subjective (much more common, arising from auditory system dysfunction).

\section*{Etiology}
\begin{itemize}
  \item Noise-induced hearing loss
  \item Age-related hearing loss (presbycusis)
  \item Meniere disease
  \item Otosclerosis
  \item Cerumen impaction (earwax)
  \item Ototoxic medications (aspirin, NSAIDs, aminoglycosides, loop diuretics)
  \item Vascular disorders (pulsatile tinnitus — AV malformations, carotid stenosis)
  \item Temporomandibular joint (TMJ) disorders
  \item Head or neck injury
  \item Acoustic neuroma (vestibular schwannoma)
\end{itemize}

\section*{Indications for Evaluation}
Seek care if:
\begin{itemize}
  \item Severe or worsening symptoms
  \item Sudden onset tinnitus, pulsatile tinnitus, tinnitus with hearing loss, dizziness, or focal neurological symptoms
  \item Accompanied by other concerning signs
\end{itemize}

\section*{Related Symptoms}
\begin{itemize}
  \item Hearing loss
  \item Vertigo
  \item Ear fullness
  \item Headache
  \item Dizziness
\end{itemize}
\textbf{Note:} Always consider serious underlying conditions when evaluating persistent tinnitus (ringing in the ears).

\section*{Self-Care / Home Management}
\begin{itemize}
  \item Rest and avoid activities that may aggravate the condition.
  \item Maintain adequate hydration and a balanced diet.
  \item Over-the-counter remedies may provide symptomatic relief (consult a pharmacist or doctor).
  \item Monitor symptoms and seek professional advice if they persist or worsen.
  \item Avoid self-medicating with prescription medications without medical guidance.
\end{itemize}

\section*{Diagnostic Approach}
\begin{itemize}
  \item A thorough history and physical examination are the first steps in evaluation.
  \item Laboratory tests (blood work, urinalysis) may be ordered based on clinical suspicion.
  \item Imaging studies (X-ray, ultrasound, CT, MRI) may be indicated when structural pathology is suspected.
  \item Specialist referral is arranged when the diagnosis remains uncertain or requires specific expertise.
\end{itemize}

\section*{Treatment / Management Overview}
\begin{itemize}
  \item Treatment targets the underlying cause identified through diagnostic evaluation.
  \item Supportive care and symptom management are provided while awaiting definitive therapy.
  \item Medications, lifestyle modifications, and physical therapies may be prescribed as appropriate.
  \item Follow-up ensures treatment effectiveness and allows timely adjustments to the care plan.
\end{itemize}

\clearpage